% -----------------------------------------------------------
\section{Reason, Reasoning}
% -----------------------------------------------------------
	\label{REASON}
	
% % ----------------------
% \subsection{See also}
% % ----------------------

% ----------------------
\subsection{1828 American Dictionary of the English Language} % *1828
% ----------------------
	\label{REASON-1828}
	
	% -----
	\subsubsection{Reason (noun)}
	% -----

		\begin{compactitem}
			\item[1] That which is thought or which is alleged in words, as the ground or cause of opinion, conclusion or determination. I have \textit{reasons} which I may choose not to disclose. You ask me my \textit{reasons}. I freely give my \textit{reasons}. The judge assigns good \textit{reasons} for his opinion, \textit{reasons} which justify his decision. Hence in general,
			\item[2] The cause, ground, principle or motive of any thing said or done; that which supports or justifies a determination, plan or measure.
			\item[3] Efficient cause. He is detained by \textit{reason} of sickness.
			\item[4] Final cause.
			\item \textit{reason}, in the English language, is sometimes taken for true and clear principles; sometimes for clear and fair deductions; sometimes for the cause, particularly the final cause.
			\item[5] A faculty of the mind by which it distinguishes truth from falsehood, and good from evil, and which enables the possessor to deduce inferences from facts or from propositions.
			\item[6] Ratiocination; the exercise of \textit{reason}.
			\item[7] Right; justice; that which is dictated or supported by \textit{reason} Every man claims to have \textit{reason} on his side.
			\item[8] \textit{reason}able claim; justice.
			\item[9] Rationale; just account.
			\item[10] Moderation; moderate demands; claims which \textit{reason} and justice admit or prescribe.
		\end{compactitem}
		
	% -----
	\subsubsection{Reason (intransitive verb)}
	% -----
	
		\begin{compactitem}
			\item[1] To exercise the faculty of \textit{reason}; to deduce inferences justly from premises. Brutes do not \textit{reason}; children \textit{reason} imperfectly.
			\item[2] To argue; to infer conclusions from premises, or to deduce new or unknown propositions from previous propositions which are known or evident. To \textit{reason} justly is to infer from propositions which are known, admitted or evident, the conclusions which are natural, or which necessarily result from them. Men may \textit{reason} within themselves; they may \textit{reason} before a court or legislature; they may \textit{reason} wrong as well as right.
			\item[3] To debate; to confer or inquire by discussion or mutual communication of thoughts, arguments or \textit{reasons}.
			\item[4] To \textit{reason} with, to argue with; to endeavor to inform, convince or persuade by argument. \textit{Reason} with a profligate son, and if possible, persuade him of his errors.
			\item[5] To discourse; to talk; to take or give an account.
		\end{compactitem}
		
	% -----
	\subsubsection{Reason (transitive verb)}
	% -----
	
		\begin{compactitem}
			\item[1] To examine or discuss by arguments; to debate or discuss. I \textit{reasoned} the matter with my friend.
			\item[2] To persuade by \textit{reasoning} or argument; as, to \textit{reason} one into a belief of truth; to \textit{reason} one out of his plan; to \textit{reason} down a passion.
		\end{compactitem}

% % ----------------------
% \subsection{Oxford English Dictionary} % *OED
% % ----------------------
	% \label{REASON-OED}

	% \begin{compactitem}
		% \item[]
	% \end{compactitem}
	
% % ----------------------
% \subsection{Buck's Theological Dictionary (1833)} % *BTD
% % ----------------------
	% \label{REASON-BTD}

	% \begin{compactdesc}
		% \item[PAGE\#]
	% \end{compactdesc}
	
% % ----------------------
% \subsection{Restoration Lexicon} % *RL *RESTORATION LEXICON
% % ----------------------
	% \label{REASON-OED}

	% \begin{compactitem}
		% \item[]
	% \end{compactitem}

% ----------------------
\subsection{Scriptural Usage} % *SCRIPTURES
% ----------------------
	\label{REASON-SCRIPTURES}

	% -----
	\subsubsection{Old Testament} % *OT
	% -----
		\label{REASON-OT}
		
		% --
		\paragraph{Reason}
		% --
	
			%
			\subparagraph{KJV}
			%
				\label{REASON-OT-KJV}
				
				I'm excluding all instances of ``by reason'' or ``by reason of.''
		
				\begin{tabular}{ll}
					Total occurrences in the OT & 11
				\end{tabular}
				
				\begin{compactdesc}
					\item[{\hyperref[1SAMUEL-12-7]{1 Samuel 12:7}}]
					\item[{\hyperref[1KINGS-9-15]{1 Kings 9:15}}]
					\item[{\hyperref[JOB-9-14]{Job 9:14}}]
					\item[{\hyperref[JOB-13-3]{Job 13:3}}]
					\item[{\hyperref[JOB-15-3]{Job 15:3}}]
					\item[{\hyperref[JOB-32-11]{Job 32:11}}]
					\item[{\hyperref[PROVERBS-26-16]{Proverbs 26:16}}] 
					\item[{\hyperref[ECCLESIASTES-7-25]{Ecclesiastes 7:25}}] 
					\item[{\hyperref[ISAIAH-1-18]{Isaiah 1:18}}]
					\item[{\hyperref[ISAIAH-41-21]{Isaiah 41:21}}] 
					\item[{\hyperref[DANIEL-4-36]{Daniel 4:36}}] 
				\end{compactdesc}
				
		% --
		\paragraph{Reasoning}
		% --
	
			%
			\subparagraph{KJV}
			%
		
				\begin{tabular}{ll}
					Total occurrences in the OT & 1
				\end{tabular}
				
				\begin{compactdesc}
					\item[{\hyperref[JOB-13-6]{Job 13:6}}]
				\end{compactdesc}
			
		% % --
		% \paragraph{Hebrew Bible}
		% % --
			% \label{REASON-HEBREW}
		
			% %
			% \subparagraph{\texorpdfstring{\he{sample word}}{}} % paste actual hebrew text here
			% %
				% \label{PAGA} % whatever is in step bible without periods
			
				% \begin{compactdesc}
					% \item[HALOT , PDF ] \textbf{} % Use the 1994 PDF
						% \begin{compactitem}
							% \item[1]
						% \end{compactitem}
					% \item[BDB , PDF ] \textbf{}
						% \begin{compactitem}
							% \item[1]
						% \end{compactitem}
					% \item[DCH PDF ] \textbf{} % don't include the actual page number, they repeat from volume to volume
						% \begin{compactitem}
							% \item[1]
						% \end{compactitem}
				% \end{compactdesc}
				
				% \begin{compactdesc}
					% \item[Some OT uses] \textbf{}
						% \begin{compactdesc}
							% \item[]
						% \end{compactdesc}
				% \end{compactdesc}
			
		% % --
		% \paragraph{Septuagint}
		% % --
			% \label{REASON-LXX}
		
			% %
			% \subparagraph{\texorpdfstring{\gr{sample word}}{}}
			% %
		
	% % -----
	% \subsubsection{New Testament} % *NT
	% % -----
		% \label{REASON-NT}
	
		% % --
		% \paragraph{KJV}
		% % --
			% \label{REASON-NT-KJV}		
	
			% \begin{tabular}{ll}
				% Total occurrences in the NT & 
			% \end{tabular}
			
			% \begin{compactdesc}
				% \item[]
			% \end{compactdesc}
			
		% % --
		% \paragraph{Greek New Testament} % *GREEK
		% % --
			% \label{REASON-GREEK}
		
			% %
			% \subparagraph{\texorpdfstring{\gr{sample word}}{}}
			% %
				% \label{KATALLAGE}
			
				% \begin{compactdesc}
					% \item[BDAG ] \textbf{}
						% \begin{compactitem}
							% \item 
						% \end{compactitem}
					% \item[CGL , PDF ] \textbf{}
						% \begin{compactitem}
							% \item 
						% \end{compactitem}
					% \item[LSJ , PDF ] \textbf{}
						% \begin{compactitem}
							% \item 
						% \end{compactitem}
					% \item[Brill , PDF ] \textbf{}
						% \begin{compactitem}
							% \item 
						% \end{compactitem}
				% \end{compactdesc}
				
				% \begin{tabular}{ll}
					% Total occurrences in the NT & 
				% \end{tabular}
				
				% \begin{compactdesc}
					% \item[Some NT uses] \textbf{}
						% \begin{compactdesc}
							% \item[]
						% \end{compactdesc}
				% \end{compactdesc}
				
			% %
			% \subparagraph{TDNT: \texorpdfstring{\gr{sample word}}{}  (3:536--556)}
			% %
				% \label{KATALLAGE-TDNT}
		
	% % -----
	% \subsubsection{Book of Mormon} % *BOM
	% % -----
		% \label{REASON-BOM}
	
		% \begin{tabular}{ll}
			% Total occurrences in the BOM & 
		% \end{tabular}
		
		% \begin{compactdesc}
			% \item[]
		% \end{compactdesc}
		
	% -----
	\subsubsection{Doctrine and Covenants} % *DC
	% -----
		\label{REASON-DC}
		
		% --
		\paragraph{Reason}
		% --
	
			\begin{tabular}{ll}
				Total occurrences in the DC & 13
			\end{tabular}
			
			\begin{compactdesc}
				\item[{\hyperref[DC3-14]{DC 3:14}}]
				\item[{\hyperref[DC45-10]{DC 45:10}}]
				\item[{\hyperref[DC45-15]{DC 45:15}}]
				\item[{\hyperref[DC49-4]{DC 49:4}}]
				\item[{\hyperref[DC50-10]{DC 50:10}}]
					\begin{compactitem}
						\item In this verse and others that talk about ``reasoning together,'' it makes me think of how Christ was interacting with the educated people in \hyperref[LUKE-2-46]{Luke 2:46--47}.
					\end{compactitem}
				\item[{\hyperref[DC50-11]{DC 50:11}}] \textbf{(2)}
				\item[{\hyperref[DC50-12]{DC 50:12}}] \textbf{(3)}
				\item[{\hyperref[DC61-13]{DC 61:13}}]
				\item[{\hyperref[DC67-3]{DC 67:3}}]
				\item[{\hyperref[DC71-8]{DC 71:8}}]
			\end{compactdesc}
			
		% --
		\paragraph{Reasoning}
		% --
		
			Note that, of the four uses below, only the DC 45 one uses ``reasoning'' as a noun.
	
			\begin{tabular}{ll}
				Total occurrences in the DC & 4
			\end{tabular}
			
			\begin{compactdesc}
				\item[{\hyperref[DC45-10]{DC 45:10}}]
				\item[{\hyperref[DC66-7]{DC 66:7}}]
					\begin{compactitem}
						\item Verb form, not a noun.
					\end{compactitem}
				\item[{\hyperref[DC68-1]{DC 68:1}}]
					\begin{compactitem}
						\item Verb form, not a noun.
					\end{compactitem}
				\item[{\hyperref[DC133-57]{DC 133:57}}]
					\begin{compactitem}
						\item Verb form, not a noun.
					\end{compactitem}
			\end{compactdesc}
		
	% % -----
	% \subsubsection{Pearl of Great Price} % *PGP
	% % -----
		% \label{REASON-PGP}
	
		% \begin{tabular}{ll}
			% Total occurrences in the PGP & 
		% \end{tabular}
		
		% \begin{compactdesc}
			% \item[]
		% \end{compactdesc}
		
% % ----------------------
% \subsection{Key Phrases} % *KEYPHRASES *PHRASES
% % ----------------------
	% \label{REASON-PHRASES}

	% \begin{compactdesc}
		% \item[] \textbf{\#times} | list
			% % \begin{compactdesc}
				% % \item[Bible anchor verses] \phantom{.}
					% % \begin{compactdesc}
						% % \item[Romans 5:5] \gr{}
					% % \end{compactdesc}
				% % \item[Commentary] ---
			% % \end{compactdesc}
	% \end{compactdesc}
	
% % ----------------------
% \subsection{Bible Dictionaries} % *DICTIONARIES
% % ----------------------
	% \label{REASON-BD}

% % ----------------------
% \subsection{Restoration Usage} % *RESTORATION
% % ----------------------
	% \label{REASON-RESTORATION}

	% % -----
	% \subsubsection{Joseph Smith} % *JS
	% % -----
		% \label{REASON-JS}
	
		% \begin{compactdesc}
			% \item[{\hyperref[KEY]{Date/Source}}]
		% \end{compactdesc}
		
	% % -----
	% \subsubsection{Temple Ordinances}
	% % -----
		% \label{REASON-TEMPLE}
	
		% \begin{compactdesc}
			% \item[{\hyperref[KEY]{Date/Source}}]
		% \end{compactdesc}
	
	% % -----
	% \subsubsection{Brigham Young} % *BY
	% % -----
		% \label{REASON-BY}
	
		% \begin{compactdesc}
			% \item[{\hyperref[KEY]{Date/Source}}]
		% \end{compactdesc}
	
% % ----------------------
% \subsection{Commentary and Studies} % *COMMENTARY *STUDIES
% % ----------------------
	% \label{REASON-COMMENTARY}

	% % -----
	% \subsubsection{Key Points}
	% % -----
		% \label{REASON-KEYS}
	
		% \begin{compactitem}
			% \item
		% \end{compactitem}
		
	% % -----
	% \subsubsection{Sample Study}
	% % -----
		% \label{REASON-study}